% !Mode:: "TeX:UTF-8"

% 中英文摘要
\begin{cabstract}

本论文围绕PyTorch深度学习框架模糊测试结果的管理与分析问题展开研究。
模糊测试可以生成大量包含边界参数、算子组合和崩溃日志的测试用例，
但这些数据在后续使用中仍存在测试产物统一管理不便、
用例价值难以直观分析、
崩溃故障线索和边界信息不易整合等问题。
针对上述问题，论文设计并实现了一套结合语义检索和大语言模型（LLM）的
模糊测试数据管理与分析系统。论文主要工作如下：

1. 针对模糊测试产物类型较多、统一管理不便的问题，
设计并实现了测试数据管理功能。
系统对测试脚本结构、执行状态、覆盖贡献数据、崩溃日志和语义向量进行统一关联，
支持测试用例导入、列表查询、条件筛选、详情查看和相似用例召回，
并通过结构化元数据和语义向量的协同存储维护测试用例之间的关联关系，
为后续数据分析和故障分析提供基础。

2. 针对用例价值分布和补充测试方向难以直观分析的问题，
设计并实现了数据分析和用例效能评估功能。
系统基于覆盖效能指数（Coverage Efficiency Index, CEI）模型，
综合考虑单个测试用例的分支覆盖贡献和加权计算成本，
并通过动态百分位阈值筛选低收益用例。
系统还在前端展示覆盖缺口、测试用例多样性分布和用例效能排行，
并提供低效用例筛选和候选补盲测试脚本生成功能，
用于辅助用户了解测试集覆盖情况和用例价值分布。

3. 针对崩溃用例中故障线索和候选边界信息不易整合分析的问题，
设计并实现了故障分析功能。
系统围绕单个崩溃用例提供故障辅助定位、相似用例查看、
候选边界条件提示和边界对比展示等功能。
同时，系统基于检索增强生成（Retrieval-Augmented Generation, RAG）方法，
整理相似用例、运行信息和候选边界信息，
生成辅助诊断报告和候选补盲测试脚本，
为用户复查崩溃原因和补充测试提供参考。

\end{cabstract}

\begin{eabstract}

This thesis studies the management and analysis of fuzzing outputs
for the PyTorch deep learning framework.
Fuzz testing can generate many test cases with boundary parameters,
unusual operator combinations along with their crash logs.
In practice, however, the resulting artifacts cannot be reused directly
in later analysis: the heterogeneous outputs are hard to keep under one schema,
the value of each case is not obvious to a reader,
and crash clues and boundary information are difficult to integrate.
To address these problems,
this thesis designs and implements a fuzzing data management and analysis system
combined with semantic retrieval and large language models (LLMs).
The main work is as follows:

1. For data management,
the thesis builds a dedicated module that handles the heterogeneous nature
of fuzzing artifacts and brings them under a single workflow.
The system associates test script structures, execution status,
coverage contribution data, crash logs, and semantic vectors in a unified way.
It supports test case import, list query, conditional filtering,
detail view, and similar-case retrieval,
and maintains relations between test cases through the coordinated storage
of structured metadata and semantic vectors,
providing a basis for later data analysis and fault analysis.

2. For data analysis,
data analysis and case efficiency evaluation functions are designed and implemented.
The system focuses on case value distribution
as well as where additional testing should be added.
A Coverage Efficiency Index (CEI) model is introduced.
For each case the model balances the branch coverage gained
against the weighted computational cost of running it,
so that low-benefit cases can be filtered out with a dynamic percentile threshold.
On the front end the system displays coverage gaps and the distribution of
case diversity, together with rankings of case efficiency.
Filtering of low-benefit cases and drafting of gap-filling test scripts
are exposed as separate operations,
giving users a clearer picture of coverage status
and of how case value is distributed across the suite.

3. For fault analysis,
a fault analysis function is designed and implemented
to integrate crash clues and candidate boundary information in failed cases.
For each crashed case, users can request localization assistance,
browse similar historical cases, inspect candidate boundary-condition hints,
or compare boundary parameters side by side.
It also uses Retrieval-Augmented Generation (RAG)
to organize similar cases, runtime information, and candidate boundary information,
and generates auxiliary diagnostic reports
and candidate gap-filling test scripts.
These functions provide references for users to review crash causes
and design additional tests.

\end{eabstract}